\section{A selective Gaussian comparison at every prior scale}
\label{sec:selective-standalone}
This section proves the free-energy comparison used above. The input is $\xi=(1,x)/\sqrt D$ with $D=d+1$, and the Gaussian precision is $\omega=N^\gamma$, $\gamma\ge1$.

\begin{lemma}[Gaussian perturbation of a sparse implementation]
Let a mean ReLU network $w_*$ have at most $m\le N$ active coordinates in every hidden layer and weighted energy $Q_N(w_*)\le E$. For $0<r\le1$, center a product Gaussian at $w_*$, retain the prior variances on coordinates entering inactive hidden rows, and multiply the prior standard deviations by $r$ on all coordinates entering active hidden rows and on every readout coordinate. Denote this comparison law by $\nu_r$. Then
\begin{align}
 \frac2\omega\operatorname{KL}(\nu_r\Vert\mu)
 &\le E+2[D+1+(H-1)m]\log(1/r),\label{eq:selective-entropy}\\
 \sup_{x\in Q}\left(\mathbb E_{\nu_r}|z_w(x)-z_{w_*}(x)|^2\right)^{1/2}
 &\le rK_E,\qquad K_E=\sqrt2(H+1)(E+1)^{H/2}.
 \label{eq:selective-rms}
\end{align}
The inactive mean rows and columns are zero, as is the mean readout off the final active coordinates.
\end{lemma}
\begin{proof}
There are at most $mD+(H-1)mN+N$ narrowed Gaussian coordinates. Since $N/\omega\le1$ and $m\le N$, their count divided by $\omega$ is at most $D+1+(H-1)m$. The Gaussian relative-entropy identity, with $r^2-1\le0$, gives \eqref{eq:selective-entropy} exactly.

Fix an input. Put $S_\ell=\|h^\ell\|_2^2/N$. The normalized first mean squared norm is at most $E$, while the first-layer noise contributes at most one to its expectation. Later mean matrices have operator norm at most $\sqrt E$. Conditional on the preceding layer, each internal noise matrix contributes at most $(N/\omega)S_{\ell-1}\le S_{\ell-1}$ to normalized expected squared norm. Centering removes cross terms and $\operatorname{ReLU}(t)^2\le t^2$, so
\[
 \mathbb E S_\ell\le2(E+1)^\ell.
\]

Let $P_\ell$ project to the active coordinates, and let
$\Delta_\ell=(\mathbb E\|P_\ell h^\ell-h_*^\ell\|_2^2/N)^{1/2}$.
The mean state is supported on those coordinates. Current-layer independence, the narrowed active-row variances, and Lipschitzness give
\[
 \Delta_1\le\sqrt2r,\qquad
 \Delta_\ell\le\sqrt E\,\Delta_{\ell-1}
       +r\sqrt{m/\omega}\sqrt{\mathbb E S_{\ell-1}}
 \le\sqrt E\,\Delta_{\ell-1}
       +\sqrt2r(E+1)^{(\ell-1)/2}.
\]
Therefore $\Delta_\ell\le\sqrt2\ell r(E+1)^{(\ell-1)/2}$. The mean readout has $\sqrt N\|v_*\|_2\le\sqrt E$; its independent perturbation contributes output variance $(r^2/\omega)S_H$. The output RMS error is thus at most
\[
 \sqrt E\,\Delta_H+r\sqrt{\mathbb E S_H/\omega}
 \le\sqrt2(H+1)r(E+1)^{H/2},
\]
proving \eqref{eq:selective-rms}.
\end{proof}

Apply the lemma to the mean implementation in \eqref{eq:path-to-Gaussian-energy}. Let $g=s(z_{w_*})$ and choose $r=(8nK_E)^{-1}$. Saturation is 1-Lipschitz, and squared loss between predictions in $[-1,1]$ and labels in $[-1,1]$ changes by at most four times their absolute prediction difference. Therefore
\[
 \mathbb E_{\nu_r}\widehat R_S(s(z_w))
 \le\widehat R_S(g)+1/(2n).
\]
The variational formula gives, for every $\lambda\ge1/n$,
\begin{equation}
 F_S(\lambda)\le\widehat R_S(g)+\lambda A,
 \qquad
 A=E+1+2[D+1+(H-1)m]\log(8nK_E).
 \label{eq:explicit-selective-cost}
\end{equation}
In particular this proves the declared resource $\widetilde O_{d,H}(1+P+m)$ without a teacher coefficient cap and without dependence on learner width.
